%=========== ΛΥΣΗ - 75ο ΘΕΜΑ Δ ===========
\begin{thema}{Δ}
Επειδή η $f$ παρουσιάζει ολικό μέγιστο στο εσωτερικό σημείο $x=2$ και
είναι παραγωγίσιμη σε αυτό, από το θεώρημα \eng{Fermat} έχουμε
\[
f'(2)=0.
\]

Αφού η $f'$ είναι γνησίως φθίνουσα στο $(0,4)$, από το
\[
f'(2)=0
\]
προκύπτει ότι
\[
f'(x)>0,\quad x\in(0,2)
\]
και
\[
f'(x)<0,\quad x\in(2,4).
\]
Άρα η $f$ είναι γνησίως αύξουσα στο $[0,2]$ και γνησίως φθίνουσα στο
$[2,4]$.

\begin{erwthma}
\item Από τα παραπάνω έχουμε
\[
f'(2)=0.
\]
Άρα η εξίσωση $f'(x)=0$ έχει τουλάχιστον μία ρίζα στο $(0,4)$, την
\[
x=2.
\]

Επειδή η $f'$ είναι γνησίως φθίνουσα στο $(0,4)$, δεν μπορεί να πάρει
την ίδια τιμή σε δύο διαφορετικά σημεία. Άρα η ρίζα αυτή είναι μοναδική.

Επομένως η εξίσωση
\[
f'(x)=0
\]
έχει ακριβώς μία ρίζα στο $(0,4)$, την
\[
{x=2}.
\]

\item Στο διάστημα $[0,2]$ η $f$ είναι συνεχής και γνησίως αύξουσα, με
\[
f(0)=0
\quad\text{και}\quad
f(2)=3.
\]
Επειδή
\[
0<\frac32<3,
\]
από το θεώρημα \eng{Bolzano} υπάρχει
\[
\alpha\in(0,2)
\]
τέτοιο ώστε
\[
f(\alpha)=\frac32.
\]
Επειδή η $f$ είναι γνησίως αύξουσα στο $[0,2]$, η ρίζα αυτή είναι
μοναδική στο $[0,2]$.

Στο διάστημα $[2,4]$ η $f$ είναι συνεχής και γνησίως φθίνουσα, με
\[
f(2)=3
\quad\text{και}\quad
f(4)=0.
\]
Επειδή
\[
0<\frac32<3,
\]
υπάρχει
\[
\beta\in(2,4)
\]
τέτοιο ώστε
\[
f(\beta)=\frac32.
\]
Επειδή η $f$ είναι γνησίως φθίνουσα στο $[2,4]$, η ρίζα αυτή είναι
μοναδική στο $[2,4]$.

Άρα η εξίσωση
\[
f(x)=\frac32
\]
έχει ακριβώς δύο ρίζες στο $(0,4)$:
\[
{\alpha\in(0,2)\ \text{και}\ \beta\in(2,4)}.
\]

\item Επειδή το $f(2)=3$ είναι ολικό μέγιστο της $f$ στο $[0,4]$,
ισχύει
\[
f(x)\leq 3,\quad x\in[0,4].
\]
Επιπλέον η ανισότητα είναι γνήσια σε υποδιάστημα, αφού
\[
f(0)=0<3.
\]
Λόγω συνέχειας της $f$, υπάρχει διάστημα κοντά στο $0$ στο οποίο
\[
f(x)<3.
\]
Άρα
\[
\int_0^4 f(x)\d x
<
\int_0^4 3\d x
=12.
\]
Επομένως
\[
{\int_0^4 f(x)\d x<12}.
\]

\item Η $f$ είναι συνεχής στο $[0,2]$ και παραγωγίσιμη στο $(0,2)$.
Άρα, από το θεώρημα μέσης τιμής, υπάρχει
\[
\xi\in(0,2)
\]
τέτοιο ώστε
\[
f'(\xi)=\frac{f(2)-f(0)}{2-0}
=\frac{3-0}{2}
=\frac32.
\]
Επομένως υπάρχει $\xi\in(0,2)$ τέτοιο ώστε
\[
{f'(\xi)=\frac32}.
\]
\end{erwthma}
\end{thema}
